\section{Guía del cuaderno}
\label{sec:guia}

\subsection{Fuentes y modo de lectura}

Llamaremos \textbf{A} a \emph{Dirichlet L-functions and the critical line},
el documento de 55 páginas, y \textbf{B} a la nota preliminar de Robin Zhang,
de 73 páginas. Las páginas citadas corresponden a los ejemplares suministrados;
sus números impresos coinciden con las páginas del PDF. Las referencias
completas se encuentran al final del cuaderno \cite{fuenteA,fuenteB}.

B presenta los objetos de manera más progresiva y explica antecedentes
matemáticos. A resulta especialmente útil para localizar hipótesis,
normalizaciones, estimaciones uniformes y nombres de declaraciones Lean.
Los utilizaremos conjuntamente: primero la motivación y después el detalle
necesario para reconstruir un paso.

El alcance que estudiamos es la hipótesis de Riemann generalizada para
\emph{funciones $L$ de Dirichlet}. La función zeta corresponde al carácter
primitivo de conductor $1$. Las formas automorfas intervienen como objetos
auxiliares; su presencia no convierte el enunciado en una afirmación sobre
todas las funciones $L$ automorfas. Véanse A, p.~1, y B, p.~3.

En estas notas distinguimos tres estados de trabajo:
\begin{description}
  \item[Desarrollado aquí.] Se incluye el cálculo o la demostración del paso.
  \item[Resultado de la fuente.] Se identifica el enunciado que utilizamos,
    sin sustituir su prueba por una explicación informal.
  \item[Pendiente de estudio.] Falta reconstruir el argumento o cotejarlo
    con mayor detalle. Esto describe el estado del cuaderno.
\end{description}

La nota B advierte expresamente, en p.~2, que su elaboración no permite
afirmar una traducción completamente fiel de todo el código. Nuestros
cotejos entre A y B tampoco constituyen una auditoría de ese código.

\subsection{La arquitectura en seis bloques}

\begin{enumerate}[label=\textbf{\arabic*.}]
  \item \textbf{Detectar ceros con una transformada.}
    La derivada logarítmica de $L$ transforma cada cero en un polo.
    Una suma suavizada sobre potencias de primos tiene una transformada de
    Mellin que conserva esos polos. La corrección en $s=1$ elimina un término
    ajeno a los ceros buscados. Esta es la Nota~\ref{sec:sumas}.
  \item \textbf{Extraer los caracteres adecuados.}
    La ortogonalidad de caracteres de un grupo de Galois abeliano expresa
    condiciones sobre Frobenius como sumas de caracteres. Los pesos y
    promedios seleccionados conservan una combinación positiva de las
    funciones corregidas. Referencias: A, §4; B, §3.
  \item \textbf{Comparar divisibilidad y normas de periodos.}
    Los periodos automorfos normalizados dan aplicaciones entre retículos.
    La divisibilidad permite ampliar los dominios sin perder la integridad
    de sus imágenes; las normas complejas controlan la altura de las
    aplicaciones. Referencias: A, §6; B, §4.1 y §4.5.
  \item \textbf{Transferir pendientes medias.}
    Una familia de mapas debe cubrir suficientes direcciones del espacio
    destino. Las estimaciones de marcos, junto con la filtración de
    Harder--Narasimhan, convierten ese control en una desigualdad para
    pendientes medias. Referencias: A, §5; B, §4.4.
  \pagebreak
  \item \textbf{Controlar los denominadores.}
    La descomposición en autoespacios puede romper la estructura integral.
    Las congruencias de Hecke y las extensiones de representaciones de Galois
    permiten medir el coste mediante grupos cohomológicos. La
    Nota~\ref{sec:proyecciones} desarrolla el ejemplo elemental de entrada.
  \item \textbf{Elegir escalas y una ventana común.}
    Después de cancelar los términos principales, aparecen costes del orden
    de $x/m$ y de $m$, salvo potencias pequeñas de $x$. Una longitud de torre
    $m$ cercana a $\sqrt{x}$ los equilibra. Las condiciones deben cumplirse
    en un mismo intervalo. Referencias: A, §§10--11; B, §§5.5--6.
\end{enumerate}

\subsection{Notación común y primeras metas}

Reservamos $G_\chi$ para la suma corregida de un carácter y $H_\psi$ para la
combinación positiva de las potencias de $\psi$. En B esta última se llama
$\mathcal G_\psi$. Usaremos $o$ para el orden del carácter, a fin de no
confundirlo con rangos de espacios automorfos. Las letras $m$ y $p$ designan,
respectivamente, una longitud de torre o una multiplicidad según se indique,
y un primo.

\begin{center}
\small
\begin{tabularx}{\textwidth}{@{}>{\raggedright\arraybackslash}p{0.34\textwidth}YY@{}}
\toprule
\textbf{Paso que queremos reconstruir} & \textbf{A} & \textbf{B} \\
\midrule
Transformada y corrección & §2, pp.~3--7 & §2, pp.~14--19 \\
Identidad de caracteres & Prop.~4.2, p.~15 & Prop.~3.1, p.~21 \\
Marcos y pendientes & Prop.~5.2 y Teo.~5.3, pp.~18--20 & Lema~4.2 y Prop.~4.3, pp.~41--44 \\
Proyecciones y denominadores & §§6--8 & Lema~5.1, p.~53; pp.~54--56 \\
Cancelación y raíz cuadrada & Props.~10.1--10.2, pp.~47--49 & (4.5)--(4.10), pp.~48--51; Lema~6.1, pp.~65--66 \\
\bottomrule
\end{tabularx}
\end{center}

La meta inmediata es poder demostrar los cálculos de las dos notas siguientes
sin consultar el texto. El siguiente desarrollo será el criterio unilateral
de ventanas: justificar la concentración del núcleo y explicar cómo se
excluye un polo a partir de la transformada. Después estudiaremos la identidad
de caracteres y la desigualdad de marcos. Las construcciones geométricas y cohomológicas
permanecen identificadas como resultados de las fuentes hasta desarrollar
sus hipótesis y demostraciones.
